\section{Scalar control and target-component constructions}
\label{sec:critical-rate-foundations}
\label{sec:critical-rate-flattening}

We establish the scalar value and the constructions used in the proof
of Theorem~\ref{thm:critical-rate-flattening}. The critical variational
functional is the one defined in the model section, with output margin
$1/2$. All logarithms are natural.

\begin{definition}[Scalar control value]
\label{def:critical-scalar-value}
For $q>0$, a scalar path of length $m\ge0$ starts at $q_0=q$.
At step $j$, choose a probability law $\eta_j$ on $\mathbb R$ with
finite entropy relative to $\mathcal N(0,q_{j-1})$, and set
\[
 q_j=\int\tanh^2 z\,d\eta_j(z).
\]
Its total cost is
\[
 \sum_{j=1}^{m}D(\eta_j\Vert\mathcal N(0,q_{j-1}))
       +\frac1{8q_m}.
\]
Let $V(q)$ be the infimum of these costs over all finite paths, including
the length-zero path. Define
\[
 V_*:=\lim_{q\downarrow0}V(q).
\]
The next lemma establishes that this limit exists and is finite.
\end{definition}

\subsection{Properties of the scalar value}

\begin{lemma}[Scalar value and increasing controls]
\label{lem:scalar-value-properties}
The function $V$ is nonincreasing, $1<V_*<2$, and there is a
universal constant $C_0$ such that
\[
 0\le V_*-V(q)\le C_0q\qquad(0<q\le1/128).
\]
For every $0<p\le1/128$, a scalar path of cost at most $V(p)+p$ can
be chosen with strictly increasing states.
\end{lemma}

\begin{proof}
Monotonicity follows from rescaling controls. If one starts a path at a
larger variance, rescale its first preactivation control by the square
root of the variance ratio. Relative entropy from the corresponding
Gaussian is unchanged, while the next second moment increases because
$|\tanh z|$ increases with $|z|$. Repeat this at every step. The final
readout cost decreases. Taking infima proves that $V$ is nonincreasing.

Put $F(u)=\mathbb E\tanh^2(\sqrt uZ)$. The elementary inequality
$\tanh^2z\ge z^2/(1+z^2)$ and Cauchy--Schwarz give
$F(u)\ge u/(1+3u)$.
Here is an elementary uniform upper bound. For $q<1/10$, replace its
reference Gaussian by $\mathcal N(0,(1+6q)q)$, at entropy cost at most
$9q^2$. The bound $F(u)\ge u/(1+3u)$ gives
\[
 q'-q\ge\frac{q^2(3-18q)}{1+3q+18q^2}>0.
\]
The cost divided by this variance increment is at most
\[
 R_0(q)=\frac{3(1+3q+18q^2)}{1-6q},
\]
which is increasing on $[0,1/10]$. Repeatedly take full steps while the
new state remains below $1/10$. Their total cost is bounded by
\[
 \int_0^{1/10}R_0(q)\,dq=-\frac{69}{200}+\log(5/2).
\]
On the last step, continuously reduce the variance inflation to hit
$1/10$ exactly; this costs at most $9/100$. The readout then costs
$5/4$. The total is at most $199/200+\log(5/2)<2$.
For starting states at least $1/10$, the length-zero path already costs
at most $5/4$. Thus $V(q)<2$ for all $q>0$.
Since $V$ is nonincreasing, its limit at zero exists and is at most two.

For a quantitative small-state bridge, let $q\le p_0:=1/128$ and
replace its reference Gaussian by
$\mathcal N(0,(1+16q)q)$. The entropy cost is at most $64q^2$.
The first Hermite coefficient
and Gaussian integration by parts give
\[
 F(u)\ge u[1-F(u)]^2\ge u(1-2u).
\]
Consequently this control gives a new state satisfying
\[
 q+4q^2\le q'\le q+16q^2\qquad(q\le p_0).
\]
The scalar exponential potential from
Lemma~\ref{lem:critical-exponential-potential}, at $t=1$ and $a=1$,
gives entropy cost at least $2(q_j-q_{j-1})$ at each step. Summing and
including the readout shows $V(q)\ge1-2q$, so $V_*\ge1$.
In fact the inequality is strict in the limit. Let
\[
 g(q)=2q-\log\mathbb E\exp\{2\tanh^2(\sqrt qZ)\}.
\]
The strict heat-equation inequality gives $g(q)>0$ for $q>0$.
Consider a path from $q<1/100$ with cost at most two. If its final
state is below $1/10$, its readout alone costs more than $5/4$.
Otherwise consider its first crossing of $1/10$.
The preceding state must exceed $1/100$: a transition with
$q'/q\ge10$ has entropy at least
$[10-1-\log10]/2>2$. To see this, its preactivation second moment
is at least $q'$, since $\tanh^2z\le z^2$; Gaussian entropy duality
gives the lower bound $[r-1-\log r]/2$ at a second-moment ratio $r\ge1$.
The pre-crossing state therefore belongs to $[1/100,1/10]$.
Summing the strengthened step bounds
$D\ge2(q'-q)+g(q)$, and minimizing the terminal
$2q_m+1/(8q_m)$, gives total cost at least
$1-2q+\min_{[1/100,1/10]}g$.
The alternative final-state-below-$1/10$ case satisfies the same bound,
since $\min_{[1/100,1/10]}g\le g(1/100)\le1/50<1/4$.
Paths of cost greater than two satisfy a stronger bound automatically.
Letting $q\downarrow0$ proves $V_*>1$.

The same Gaussian-growth construction, stopped exactly at any
$p\le p_0$, costs at most $17p$ when started at an arbitrarily small
positive state. Indeed all full steps cost at most 16 times their
variance increment. For the last step, continuously reduce the variance
inflation to hit $p$ exactly; its cost is at most $64p^2\le p$.
Appending a nearly optimal path from $p$ and letting the starting state
tend to zero yields $V_*\le17p+V(p)$.

We also need a strict dissipation estimate. There is a universal $c>0$
such that
\begin{equation}\label{eq:strict-critical-potential}
 \log\mathbb E\exp\{\tanh^2(\sqrt qZ)\}\le q-cq^2,
 \qquad 0\le q\le1.
\end{equation}
For a direct verification, $G(z)=\exp(\tanh^2z)$ satisfies $G''<2G$
away from zero. The heat equation gives a strict inequality
$\log\mathbb EG(\sqrt qZ)<q$ for $q>0$.
Taylor expansion gives
$\log\mathbb EG(\sqrt qZ)=q-q^2+O(q^3)$ at zero. Compactness on
$[0,1]$ therefore gives a positive $c$ in~\eqref{eq:strict-critical-potential}.
Entropy duality now implies
\begin{equation}\label{eq:scalar-step-dissipation}
 D(\eta_j\Vert\mathcal N(0,q_{j-1}))
 \ge q_j-q_{j-1}+cq_{j-1}^2.
\end{equation}

An arbitrary finite path can be made strictly increasing without raising
its cost. If a step decreases or preserves the state, delete that step.
Start the remaining controls at the larger preceding state and rescale
each successive control as in the monotonicity argument. Their entropies
are unchanged and the final readout cost cannot increase. Repeat until
there are no nonincreasing steps; the procedure terminates because each
deletion shortens the path. Take a path within $p$ of $V(p)$ and apply
this procedure. The resulting increasing path has the asserted cost, proving the
remaining claim.
\end{proof}

\subsection{A small-variance layer forces the scalar cost}

\begin{lemma}[Comparison of a signed pair with the scalar control problem]
\label{lem:critical-pair-scalar-comparison}
For any admissible critical variational path and any signed input pair,
put $b=(e_x+se_y)/2$, $s\in\{-1,1\}$, and $v_j=b^TQ_jb$.
If the output satisfies $|b^Tz|\ge1/2$, the part of the path after
layer $j$, including the readout, costs at least $V(v_j)$.
Moreover, a transition from pair variance $v\le1$ to $v'$ has
entropy cost at least $v'-v+cv^2$, with the constant from
\eqref{eq:strict-critical-potential}.
\end{lemma}

\begin{proof}
It suffices to compare one transition. Lift a finite-entropy feature
law through the inverse coordinatewise tanh map to its preactivation
law. Under the reference law, the preactivation combination $b^TG$ is
$\mathcal N(0,v)$. Under the controlled law, denote its distribution by
$\eta$. Data processing bounds its Gaussian relative entropy by the
original layer entropy. The pair contraction inequality gives
\[
 \int\tanh^2z\,d\eta(z)\ge v',
\]
where $v'$ is the next feature-pair second moment. If the left side is
$q'$, the definition of $V$ and its monotonicity imply
\[
 V(v)\le D(\eta\Vert\mathcal N(0,v))+V(q')
       \le D(\nu'\Vert\mathsf M_Q)+V(v').
\]
The final readout costs at least $1/(8v_{L-1})\ge V(v_{L-1})$.
Telescoping proves the assertion. Zero-variance or singular-covariance
cases cause no problem: a finite-cost path reaching the stipulated
nonzero pair output must have positive pair variance at every relevant
step, and Cauchy--Schwarz is taken on the covariance range.
For the final assertion, entropy duality applied to
\eqref{eq:strict-critical-potential} gives
\[
 D(\eta\Vert\mathcal N(0,v))
 \ge \int\tanh^2z\,d\eta(z)-v+cv^2
 \ge v'-v+cv^2.
\]
Data processing transfers this bound to the full layer entropy.
\end{proof}

\subsection{Holding and growing an arbitrarily weak target component}

Suppose a feature second moment has the decomposition
\[
 Q=vff^T+R,\qquad R\succeq0,\qquad R(x,x)=r\quad\hbox{for every }x.
\]
Its diagonal is $q=v+r$. Inflate only the target component, taking the
next preactivation covariance to be $Q+\epsilon vff^T$.
Its entropy cost is at most $\epsilon^2/4$. Indeed, $f$ lies in the
range of $Q$, and $u=vf^TQ^+f\le1$. The Gaussian relative entropy
on that range is
$[\epsilon u-\log(1+\epsilon u)]/2\le\epsilon^2/4$;
applying tanh cannot increase it. Put
\[
 \psi_r(z)=\mathbb E\tanh(z+\sqrt rZ),\qquad
 v'=\mathbb E\psi_r(\sqrt{(1+\epsilon)v}Z)^2,
 \qquad q'=F(q+\epsilon v).
\]
Conditional expectation gives $Q'\succeq v'ff^T$.
Conditional Gaussian Poincare, using the unit Lipschitz constant of tanh,
also gives
\begin{equation}\label{eq:critical-residual-nonincreasing}
 r':=q'-v'\le r.
\end{equation}
The map $\epsilon\mapsto v'$ is continuous and nondecreasing. At
$\epsilon=0$ it is at most $v$, while the first Hermite projection gives
\begin{equation}\label{eq:critical-component-growth}
 v'\ge(1+\epsilon)v[1-F(q+\epsilon v)]^2
       \ge(1+\epsilon)v[1-2(q+\epsilon v)].
\end{equation}
These statements depend only on the tracked scalars $v,r,\epsilon$,
not on the remaining entries of $R$ or on the choice of $f$.

\begin{lemma}[Cheap holding and prescribed growth]
\label{lem:critical-cheap-holding}
There are universal constants $A,B,C_2$ with the following properties.
Let $D=d+1$ and $0<\delta\le p\le1/128$.
Starting from the unperturbed first feature law, at zero entropy cost
and after at most $B/\delta$ hidden layers, one can track
\[
 Q\succeq v_0ff^T,\qquad v_0=\delta e^{-AD},\qquad Q(x,x)\le\delta.
\]
One may hold this $v_0$ exactly fixed for any specified number $H$ of
layers, at total entropy cost at most $64H\delta^2$, retaining diagonal
residual at most $\delta$.
Thereafter one can grow the tracked component exactly to $p$ in at most
$BD/\delta$ layers, with entropy cost at most
\[
 C_2p+C_2D\delta,
\]
and with residual diagonal still at most $\delta$.
The number of growth layers is fixed in advance from $v_0,\delta,p$,
so it is independent of the preceding holding time and of $f$.
\end{lemma}

\begin{proof}
The first unperturbed covariance has smallest eigenvalue at least
$e^{-13D}$ by Lemma~\ref{lem:first-tanh-spectrum}. Since the cube has
$2^d\le e^D$ points, it dominates $e^{-14D}ff^T$ for every $f$.
Lemma~\ref{lem:critical-free-cooling} shows that, after $O(1/\delta)$
unmodified layers, the diagonal lies in $[\delta/2,\delta]$ and the
tracked target component is at least $b\delta e^{-14D}$.
Enlarging $A$ gives the displayed common
lower component $v_0$; reduce the tracked component to precisely that
value if necessary.

For holding, use an inflation $\epsilon\in[0,16\delta]$ chosen to
make $v'=v_0$. Such a choice exists by continuity: at zero inflation
$v'\le v_0$, and~\eqref{eq:critical-component-growth} gives $v'\ge v_0$
at $16\delta$ when $q\le\delta\le1/128$.
The residual does not increase, so the same argument applies at every
holding layer. Each costs at most $64\delta^2$.

For growth, prescribe the scalar sequence
\[
 v_{j+1}=\min\{p,\ v_j[1+2(\delta+v_j)]\},
\]
stopping upon reaching $p$. Throughout it, $r\le\delta$ and $v_j\le p$.
At inflation $\epsilon=16(\delta+v_j)$,
\eqref{eq:critical-component-growth} gives
$v'\ge v_j[1+4(\delta+v_j)]$.
Therefore an inflation no larger than this realizes the prescribed next
value exactly. Its entropy cost is at most $64(\delta+v_j)^2$.
For a full step, with $\Delta v=2(\delta+v_j)v_j$, this cost is
\[
 32\Delta v+32\delta\frac{\Delta v}{v_j}.
\]
Here $\Delta v/v_j\le4p<1$, so
$\Delta v/v_j\le2\log(v_{j+1}/v_j)$.
Summing, bounding the possible last step separately by $256p^2$, and
using $\delta\log(p/\delta)\le p/e$, proves a bound
$C_2p+C_2D\delta$.

Before $v$ reaches $\delta$, each full step multiplies it by at least
$1+2\delta$, so at most $AD/\delta+1$ steps are needed.
Afterwards each full step decreases $1/v$ by at least one, so at most
$1/\delta+1$ further full steps and one final step are needed.
This proves the stated $BD/\delta$ length bound. The prescribed scalar
sequence fixes the length independently of the residual or holding time.
\end{proof}

\subsection{Stable lifting of scalar controls}

\begin{lemma}[Lifting with a small residual]
\label{lem:critical-scalar-lift}
Suppose $Q\succeq pff^T$ has constant diagonal $p+r$, with
$0\le r\le\delta$. Take a scalar path of length $m$ starting at $p$.
It can be lifted to a vector feature path with no greater total hidden
entropy. Its final second moment dominates
\[
 (1-\delta)^{2m}q_mff^T.
\]
The residual diagonal stays at most $\delta$ throughout. If the scalar
path has total cost at most three and $2m\delta\le1/2$, the lifted path,
including output $z=f/2$, costs at most the scalar cost plus $12m\delta$.
\end{lemma}

\begin{proof}
For $z\ge0$ and $r\ge0$, symmetry of a standard Gaussian gives
\[
 \psi_r(z)=\mathbb E\tanh(z+\sqrt rZ)
 \ge[1-F(r)]\tanh z\ge(1-r)\tanh z.
\]
Indeed, writing $u=\tanh z$ and $w=\tanh(\sqrt rZ)$, the average of
the two values corresponding to $w$ and $-w$ is
$u(1-w^2)/(1-u^2w^2)\ge u(1-w^2)$. Oddness gives the corresponding
absolute-value inequality for all real $z$.

Suppose the current tracked vector component is $v$, whereas the
nominal scalar state is $q$. Rescale the nominal scalar preactivation
control $\eta$ by $\sqrt{v/q}$. Its entropy relative to
$\mathcal N(0,v)$ equals the nominal entropy relative to
$\mathcal N(0,q)$. Independently add a centered Gaussian with the
residual covariance. Data processing bounds the vector feature entropy
by that same scalar entropy. Its new tracked component is the second
moment of the conditional feature mean.
Since $\tanh(cz)\ge c\tanh z$ for $z\ge0$ and $0\le c\le1$,
the new component $v'$ obeys
\[
 \frac{v'}{q'}\ge(1-\delta)^2\min\{1,v/q\},
\]
where $q'$ is the nominal next scalar state.
The conditional-variance argument from
\eqref{eq:critical-residual-nonincreasing} applies equally to a
non-Gaussian scalar control, so the residual diagonal never increases.
Starting from $v=q=p$ proves the component bound by induction.

The lifted readout cost is at most the nominal readout cost divided by
$(1-\delta)^{2m}\ge1-2m\delta$. If $2m\delta\le1/2$, its increase is
at most $4m\delta$ times the nominal readout cost, which is at most three.
This proves the final assertion.
\end{proof}

